\documentclass[10pt,a4paper]{amsart}
\usepackage[margin=19mm]{geometry}
\usepackage{amsmath,amssymb,mathtools}
\usepackage[scr=rsfs,cal=euler]{mathalfa}
\usepackage{xfrac}
\usepackage[unicode,hidelinks,bookmarks=false]{hyperref}
\newcommand{\ie}{i.e.\ }
\newcommand{\cf}{cf.\ }
\newcommand{\resp}{respectively\ }
\newcommand{\Subsection}{\subsection}
\providecommand{\nobreakdash}{\nobreak\hbox{-}}
\setlength{\emergencystretch}{2em}
\multlinegap0pt
\usepackage{bm}
\usepackage{bbm}
\usepackage{dsfont}
\usepackage{xspace}
\usepackage{cancel}
\usepackage{mathtools}
\usepackage{amsthm}
\makeatletter
\@ifundefined{theorem}{\newtheorem{theorem}{Theorem}}{}
\@ifundefined{corollary}{\newtheorem{corollary}[theorem]{Corollary}}{}
\@ifundefined{lemma}{\newtheorem{lemma}[theorem]{Lemma}}{}
\@ifundefined{proposition}{\newtheorem{proposition}[theorem]{Proposition}}{}
\makeatother
\newcommand{\BR}{{\rm{Br}}}
\newcommand{\Si}{\Sigma}
\newcommand{\ZR}{\mathbb{R}}
\newcommand{\ZT}{\mathbb{T}}
\newcommand{\cb}{\mathcal B}
\newcommand{\cl}{\mathcal L}
\newcommand{\de}{\delta}
\newcommand{\e}{\varepsilon}
\newcommand{\ga}{\gamma}
\newcommand{\si}{\sigma}
\newcommand{\supp}{{\rm supp}}
\newcommand{\wt}{\widetilde}
\title[Weighted $L^2$ estimates with applications to $L^p$ problems]{Weighted $L^2$ estimates with applications to $L^p$ problems}
\author{Shukun Wu}
\date{Source: arXiv:2506.02650v1 (CC BY 4.0)}
\begin{document}
\maketitle

\noindent\emph{Source and licence.} Weighted $L^2$ estimates with applications to $L^p$ problems. Version arXiv:2506.02650v1 (CC BY 4.0), distributed under the Creative Commons Attribution 4.0 International licence, as recorded in the arXiv licence metadata for this exact version and shown on the abstract page https://arxiv.org/abs/2506.02650v1. Full rights notice: licenses/arxiv-2506.02650-CC-BY-4.0.txt.\par\smallskip

\noindent\textit{Editorial scope.} Counted unit: the Proposition whose source label is \texttt{weighted-lp-prop} in the article (source0.tex lines 1078--1086, source label \texttt{weighted-lp-prop}), delivered together with the article's own complete original proof of it (source0.tex lines 1088--1161, one \texttt{proof} environment of 74 lines). No mathematics is written, repaired or re-derived: statement and proof are the article's own text line for line. Required context retained uncounted: wpt (lines 597--605); katz-tao-set-lem (lines 646--655); main-cor (lines 1007--1018); broad-narrow-lem (lines 1050--1059); parabolic-rescaling (lines 1064--1067). Every definition, display and label of those lines is retained unchanged. Omitted from this package and not redistributed: 1--596, 606--645, 656--1006, 1019--1049, 1060--1063, 1068--1077, 1087--1087, 1162--1472. Nothing inside a retained excerpt is omitted or altered: every retained body is delivered exactly as the article's lines are written, so no omission receipt of any kind accompanies this package. Citations: the retained excerpts carry no citation command at all, so no bibliography entry of the article is transcribed and no reference is redistributed. Reference adaptation: none was needed and none was made; every reference inside the delivered unit resolves inside it, so no unresolved reference or citation is produced. Printed slips kept literal: no printed word, symbol, hypothesis or number of the counted statement or of its proof was altered. Layout and class: the article's own class is \texttt{amsart} with packages including tikz, xcolor, amssymb, latexsym, amsmath, extarrows, graphicx, mathrsfs, hyperref, url; the wrapper is the lane's pinned \texttt{amsart} editorial wrapper at 10pt on A4, which restates the theorem-like environments the retained text needs and adds the article's own macro definitions that the retained text needs (\texttt{\textbackslash BR}, \texttt{\textbackslash Si}, \texttt{\textbackslash ZR}, \texttt{\textbackslash ZT}, \texttt{\textbackslash cb}, \texttt{\textbackslash cl}, \texttt{\textbackslash de}, \texttt{\textbackslash e}, \texttt{\textbackslash ga}, \texttt{\textbackslash si}, \texttt{\textbackslash supp}, \texttt{\textbackslash wt}), with extra wrapper packages (bm, bbm, dsfont, xspace, cancel, mathtools). The delivered numbering is the unit's own. Assets: no retained span contains a figure environment, a graphics-inclusion command, a PDF-page inclusion, a table or a TikZ picture; no image file is copied into this package, the package declares no asset dependency and no image file is redistributed with it. Licence: version arXiv:2506.02650v1 of this article is distributed under the Creative Commons Attribution 4.0 International licence, as recorded in the arXiv licence metadata for this exact version and shown on the abstract page https://arxiv.org/abs/2506.02650v1. A full rights notice is delivered at licenses/arxiv-2506.02650-CC-BY-4.0.txt. Characters: every retained excerpt is pure ASCII, so the delivered engine, XeLaTeX with its default Latin Modern text font, reports no missing character.
\begin{lemma}
	\label{wpt}
	The wave packet decomposition satisfies the following properties.
	\begin{enumerate}\item $Ef=\sum_{T\in\bar\ZT}Ef_{T}$.
		\item $|Ef_{T}(x)|\lesssim R^{-1000}$ when $x\in B_R\setminus T$.
		\item $\supp f_{T}\subset 3\theta$ when $T$ has direction $V(\theta)$.
		\item $\{V(\theta)\}_{\theta\in\Theta}$ are $\gtrsim R^{-1/2}$-separated.
	\end{enumerate}    
\end{lemma}

\begin{lemma}
\label{katz-tao-set-lem}
Let $\de\in(0,1)$, and let $X\subset[0,1]^2$ be a union of $\de$-balls.
Let $\ga\geq1$ be defined as 
\begin{equation}
\label{katz-tao-constant}
    \ga=\ga_X=\sup_{r\in[\de,1]}\sup_{x\in[0,1]^2}\frac{|X\cap B(x,r)|\de^{-2}}{r/\de}.
\end{equation}
Then there exists $X'\subset X$ with $|X|\approx\ga^{-1}|X|$ such that $X'$ is a Katz-Tao $(\de,1)$-set.
\end{lemma}

\begin{corollary}
\label{main-cor}
Let $R\geq 1$ and $\e\in(0,10^{-3})$.
Let $K\leq R^{\e^4}$, $\e'\in(0,\log K/\log R)$, and let $\Si=\{\si\}$ be a collection of finitely overlapping $K^{-1}$-caps. 
Let $X\subset B_R$ be a union of unit balls such that the $R^{-1}$-dilate of  $X$ is a Katz-Tao $(R^{-1},1)$-set.
Suppose $A\geq R^{\e'}$.
Then when $q\geq 18/5$,
\begin{equation}
\label{main-esti-cor}
    \|\BR_A Ef\|_{L^q(X)}\leq C_{\e,\e'} R^{\e}\|f\|_2.
\end{equation}
\end{corollary}

\begin{lemma}[Broad-narrow]
	\label{broad-narrow-lem}
	Let $K\geq10$.
	Let $\Si$ be a family of $K^{-1}$-intervals that forms a finite-overlapping cover of $[-1,1]$.
	Then for all $x\in\ZR^2$ and $A\in[1,K]$,  
	\begin{align}
	\label{broad-narrow}
		|Ef(x)|\lesssim A\max_{\si\in\Si}|Ef_{\si}(x)|+K\cdot \BR_AEf(x).
	\end{align}
\end{lemma}

\begin{equation}
\label{parabolic-rescaling}
    \cl_\si(x_1,x_2)=(K^{-1}(x_1+2x_2\xi_\si), K^{-2}x_2).
\end{equation}

\begin{proposition}
\label{weighted-lp-prop}
Let $X\subset B_R$ be a union of unit balls that the $R^{-1}$-dilate of $X$ is a Katz-Tao $(R^{-1},1)$-set.
Then when $q\in[18/5,4]$ and $2/q+1/p=1$, for all $\e>0$,
\begin{equation}
\label{weighted-lp-esti}
    \|Ef\|_{L^q(X)}\leq C_\e R^\e\|f\|_p.
\end{equation}
\end{proposition}

\begin{proof}
For a given $\e$, pick $K=R^{\e^8}$ and $A=K^{\e^2}$.
By Lemma \ref{broad-narrow-lem}, we have
\begin{equation}
\label{broad-narrow-app-1}
    \|Ef\|_{L^q(X)}^q\lesssim K^{O(\e^2)}\sum_\si\int_X|Ef_\si|^q+K^{O(1)}\int_X|\BR_A Ef|^q.
\end{equation}
We distinguish two cases.

\smallskip

{\bf Case 1.} Suppose the first terms in \eqref{broad-narrow-app-1} dominates.
Fix $\si$.
Let $\cb_\si$ be a family of finite-overlapping $K\times K^2$-rectangles oriented with direction $\si$ that forms a cover of $X$.
By dyadic pigeonholing, there exist $s\in[1,K^3]$ and a set $\cb_{\si,s}\subset \cb_\si$ such that the following is true:
\begin{enumerate}
    \item For all $B_\si\in\cb_{\si,s}$, $|X\cap B_\si|\sim s$.
    \item
    We have, since $|Ef_\si|$ is essentially constant on each $B_\si\in\cb_{\si,s}$,
    \begin{equation}
        \int_X|Ef_\si|^q\lessapprox \int_{X\cap (\cup_{\cb_{\si,s}})}|Ef_\si|^q\lesssim sK^{-3}\int_{\cup_{\cb_{\si,s}}}|Ef_\si|^q.
    \end{equation}
\end{enumerate}
Similar to the proof of Corollary \ref{main-cor}, either we can easily conclude \eqref{weighted-lp-esti}, or by dyadic pigeonholing, there exists a union of $K\times K^2$-rectangles $\wt X_\si\subset\cup_{\cb_{\si,s}}$ so that 
\begin{enumerate}
    \item $\|Ef_\si\|_{L^q(B_\si)}$ are about the same for all $K\times K^2$-rectangles $B_\si\subset \wt X_\si$.
    \item We have
    \begin{equation}
    \label{uniform-in-K-rect}
        \int_X|Ef_\si|^q\lessapprox sK^{-3}\int_{\wt X_\si}|Ef_\si|^q.
    \end{equation}
\end{enumerate}


For any $r\in[1,R/K^2]$, let $Q_r$ be an $rK\times rK^2$-rectangle oriented with direction $\si$.
Since the $R^{-1}$-dilate of $X$ is a Katz-Tao $(R^{-1},1)$-set, we know that $|X\cap Q_r|=\sum_{B_r\subset Q_r}|X\cap B_r|\lesssim K^2 r$, where $\{B_r\}$ are finite-overlapping $rK$-balls.
As a result
\begin{equation}
    \#\{B_\si\in\cb_{\si,s}:B_\si\subset Q_r\}\lesssim (K^2s^{-1})r.
\end{equation}
By Lemma \ref{katz-tao-set-lem}, there exists a union of $K\times K^2$-rectanlges $\wt X_\si'\subset \wt X_\si$ such that
\begin{enumerate}
    \item $|\wt X_\si'|\gtrapprox s K^{-2}|\wt X_\si|$.
    \item Recall \eqref{parabolic-rescaling}.
    Inside an unit ball, the $(R/K^2)^{-1}$-dilate of $\cl_\si(\wt X_\si')$ is a Katz-Tao $((R/K^2)^{-1},1)$-set.
\end{enumerate}
Let $g(\xi)=f(K^{-1}(\xi-\xi_\si))$, where $\xi_\si$ is the center of $\si$.
Put the above information back to \eqref{uniform-in-K-rect} so that
\begin{equation}
\label{after-rescaling}
    \int_X|Ef_\si|^q\lessapprox K^{-1}\int_{\wt X_\si'}|Ef_\si|^q=K^{2-q}\int_{\cl_\si(\wt X_\si')}|Eg|^q.
\end{equation}

Note that $2/q+1/p=1$.
Apply \eqref{weighted-lp-esti} at scale $R/K^2$ and use Lemma \ref{wpt} to get
\begin{equation}
\label{apply-recaling}
    \int_{\cl_\si(\wt X_\si')}|Eg|^q\leq C_\e (R/K^2)^{\e}\|g\|_p^q=K^{-2\e}K^{q-2}C_\e R^\e\|f_\si\|_p^q.
\end{equation}
Put \eqref{apply-recaling} and \eqref{after-rescaling} back to \eqref{broad-narrow-app-1} and sum up $\si$ to obtain (note that $p\leq q$)
\begin{equation}
    \|Ef\|_{L^q(X)}^q\lessapprox K^{O(\e^2)-2\e}C_\e R^\e\sum_\si\|f_\si\|_p^q\leq C_\e R^\e\|f\|_p^q.
\end{equation}


\smallskip

{\bf Case 2.} Suppose the second terms in \eqref{broad-narrow-app-1} dominates.
Then, by Corollary \ref{main-cor} with $\e:=\e^2$, we have (note that $p\geq 2$ when $q\in[18/5,4]$)
\begin{equation}
    \|Ef\|_{L^q(X)}^q\lesssim K^{O(1)}\|\BR_A Ef\|_{L^q(X)}^q\lesssim R^{O(\e^2)}\|f\|_2^q\lesssim R^{O(\e^2)}\|f\|_p^q. \qedhere
\end{equation}

\end{proof}

\end{document}
